% Auto-generated by scripts/generate_blueprint.py
% Do not edit manually.

\chapter{External Dependencies}

\begin{definition}[Definition of a field and basic field theory (characteristic, prime fields, field extensions)]
\label{ext:field}
\notready
Category: undergraduate\_prerequisite.
\end{definition}

\begin{definition}[Definition of a commutative ring with identity, ring homomorphisms, ring extensions]
\label{ext:ring}
\notready
Category: undergraduate\_prerequisite.
\end{definition}

\begin{definition}[Definition of an integral domain]
\label{ext:integral-domain}
\notready
Category: undergraduate\_prerequisite.
\end{definition}

\begin{definition}[Ideals (prime, maximal, principal), quotient rings, ideal generation]
\label{ext:ideal-theory}
\notready
Category: undergraduate\_prerequisite.
\end{definition}

\begin{definition}[The integer ring Z and its basic properties]
\label{ext:integers}
\notready
Category: undergraduate\_prerequisite.
\end{definition}

\begin{definition}[The field Q of rational numbers as fraction field of Z]
\label{ext:rationals}
\notready
Category: undergraduate\_prerequisite.
\end{definition}

\begin{definition}[The real numbers R, R\_\{>=0\}, R\_\{>0\}, and the standard ordering on R]
\label{ext:real-numbers}
\notready
Category: undergraduate\_prerequisite.
\end{definition}

\begin{definition}[Prime numbers, divisibility in Z, coprime integers]
\label{ext:primes-and-divisibility}
\notready
Category: undergraduate\_prerequisite.
\end{definition}

\begin{definition}[Every nonzero integer has a unique prime factorization (up to units)]
\label{ext:fundamental-theorem-of-arithmetic}
\notready
Category: undergraduate\_prerequisite.
\end{definition}

\begin{definition}[Finite fields F\_p and F\_q, existence and uniqueness up to isomorphism]
\label{ext:finite-fields}
\notready
Category: undergraduate\_prerequisite.
\end{definition}

\begin{definition}[Polynomial rings A[x], monic polynomials, irreducible polynomials, minimal polynomials]
\label{ext:polynomial-ring}
\notready
Category: undergraduate\_prerequisite.
\end{definition}

\begin{definition}[Fraction field (field of fractions) of an integral domain]
\label{ext:fraction-field}
\notready
Category: undergraduate\_prerequisite.
\end{definition}

\begin{definition}[Principal ideal domains: every ideal is generated by a single element]
\label{ext:pid}
\notready
Category: undergraduate\_prerequisite.
\end{definition}

\begin{definition}[Unique factorization domains]
\label{ext:ufd}
\notready
Category: undergraduate\_prerequisite.
\end{definition}

\begin{definition}[Noetherian rings: every ideal is finitely generated, ascending chain condition]
\label{ext:noetherian-ring}
\notready
Category: undergraduate\_prerequisite.
\end{definition}

\begin{definition}[Krull dimension of a ring (supremum of lengths of chains of prime ideals)]
\label{ext:krull-dimension}
\notready
Category: undergraduate\_prerequisite.
\end{definition}

\begin{definition}[Localization of a ring at a prime ideal]
\label{ext:localization}
\notready
Category: undergraduate\_prerequisite.
\end{definition}

\begin{definition}[Group homomorphisms (used for valuations as homomorphisms k* -> R)]
\label{ext:group-homomorphism}
\notready
Category: undergraduate\_prerequisite.
\end{definition}

\begin{definition}[Power series ring k[[t]] and Laurent series field k((t))]
\label{ext:power-series-and-laurent-series}
\notready
Category: undergraduate\_prerequisite.
\end{definition}

\begin{definition}[Existence of algebraic closure of a field, field embeddings]
\label{ext:algebraic-closure}
\notready
Category: undergraduate\_prerequisite.
\end{definition}

\begin{definition}[Finite field extensions, algebraic extensions, degree of extension]
\label{ext:field-extensions}
\notready
Category: undergraduate\_prerequisite.
\end{definition}

\begin{definition}[Symmetric functions of roots: coefficients of a polynomial are symmetric functions of its roots]
\label{ext:symmetric-polynomials}
\notready
Category: undergraduate\_prerequisite.
\end{definition}

\begin{definition}[Free Z-modules, Z-lattices in a Q-vector space]
\label{ext:z-module}
\notready
Category: undergraduate\_prerequisite.
\end{definition}

\begin{definition}[Regular local rings (dim m/m\textasciicircum{}2 = dim A)]
\label{ext:regular-local-ring}
\notready
Category: undergraduate\_prerequisite.
\end{definition}

\begin{definition}[Vector spaces over a field, dimension]
\label{ext:vector-space}
\notready
Category: undergraduate\_prerequisite.
\end{definition}

\begin{definition}[The triangle inequality for real-valued functions]
\label{ext:triangle-inequality}
\notready
Category: undergraduate\_prerequisite.
\end{definition}

\begin{definition}[If r/s in lowest terms is a root of a monic polynomial in Z[x], then s = +-1]
\label{ext:rational-root-test}
\notready
Category: external\_result.
\end{definition}

\begin{definition}[Every discrete subgroup of R is isomorphic to Z]
\label{ext:discrete-subgroups-of-R}
\notready
Category: folklore.
\end{definition}

\begin{definition}[Z is a principal ideal domain]
\label{ext:Z-is-PID}
\notready
Category: folklore.
\end{definition}

\begin{definition}[F\_q[t] is a principal ideal domain]
\label{ext:Fq-t-is-PID}
\notready
Category: folklore.
\end{definition}

\begin{definition}[Every field contains exactly one prime field (Q or F\_p) according to its characteristic]
\label{ext:every-field-contains-prime-field}
\notready
Category: folklore.
\end{definition}

\begin{definition}[Z has Krull dimension one]
\label{ext:Z-krull-dim-one}
\notready
Category: folklore.
\end{definition}

\begin{definition}[Riemann zeta function: Euler product, analytic continuation, simple pole at s=1]
\label{ext:riemann-zeta}
\notready
Category: external\_result.
\end{definition}

\begin{definition}[Proof of equivalent characterizations of DVRs (references [1, §23] or [2, §9])]
\label{ext:dvr-characterizations-proof}
\notready
Category: external\_result.
\end{definition}

\begin{definition}[Proof that integrality is transitive in ring extension towers (references [1, Thm. 10.27] or [2, Cor. 5.4])]
\label{ext:transitivity-of-integrality-proof}
\notready
Category: external\_result.
\end{definition}

\chapter{Lecture 1}

\begin{remark}[Introduction to number theory]
\label{Lecture1:Introduction}
\uses{ext:field, ext:ring, ext:ideal-theory, ext:integers, ext:rationals, ext:primes-and-divisibility, ext:fundamental-theorem-of-arithmetic, ext:finite-fields, ext:polynomial-ring, ext:fraction-field, ext:pid, ext:krull-dimension, ext:field-extensions, ext:Z-is-PID, ext:Fq-t-is-PID, ext:every-field-contains-prime-field, ext:Z-krull-dim-one, ext:riemann-zeta}
\discussion{Lecture1:Introduction}
[Introduction: Introduction to number theory]
\end{remark}

\begin{remark}[Rings have multiplicative identity]
\label{Lecture1:Remark1.1}
\uses{Lecture1:Introduction, ext:ring}
\discussion{Lecture1:Remark1.1}
[Remark: Rings have multiplicative identity]
\end{remark}

\begin{remark}[Section 1.2: Absolute values]
\label{Lecture1:Discussion_after_1.1}
\uses{Lecture1:Remark1.1, ext:field}
\discussion{Lecture1:Discussion_after_1.1}
[Discussion: Section 1.2: Absolute values]
\end{remark}

\begin{definition}[Absolute value on a field]
\label{Lecture1:Definition1.2}
\uses{Lecture1:Discussion_after_1.1, ext:field, ext:real-numbers, ext:triangle-inequality}
[Definition: Absolute value on a field]
\end{definition}

\begin{example}[Trivial absolute value]
\label{Lecture1:Example1.3}
\uses{Lecture1:Definition1.2, ext:field, ext:real-numbers}
[Example: Trivial absolute value]
\end{example}

\begin{lemma}[Nonarchimedean characterization]
\label{Lecture1:Lemma1.4}
\uses{Lecture1:Example1.3, ext:field}
[Lemma: Nonarchimedean characterization]
\end{lemma}

\begin{corollary}[Positive characteristic implies nonarchimedean]
\label{Lecture1:Corollary1.5}
\uses{Lecture1:Lemma1.4, ext:field, ext:finite-fields}
[Corollary: Positive characteristic implies nonarchimedean]
\end{corollary}

\begin{definition}[Equivalent absolute values]
\label{Lecture1:Definition1.6}
\uses{Lecture1:Corollary1.5, ext:field, ext:real-numbers}
[Definition: Equivalent absolute values]
\end{definition}

\begin{remark}[Section 1.3: Absolute values on Q]
\label{Lecture1:Discussion_after_1.6}
\uses{Lecture1:Definition1.6, ext:field, ext:integers, ext:rationals, ext:primes-and-divisibility, ext:fundamental-theorem-of-arithmetic}
\discussion{Lecture1:Discussion_after_1.6}
[Discussion: Section 1.3: Absolute values on Q]
\end{remark}

\begin{definition}[p-adic valuation and absolute value]
\label{Lecture1:Definition1.7}
\uses{Lecture1:Discussion_after_1.6, ext:integers, ext:rationals, ext:primes-and-divisibility, ext:fundamental-theorem-of-arithmetic}
[Definition: p-adic valuation and absolute value]
\end{definition}

\begin{theorem}[Ostrowski's Theorem]
\label{Lecture1:Theorem1.8}
\uses{Lecture1:Definition1.7, ext:rationals}
[Theorem: Ostrowski's Theorem]
\end{theorem}

\begin{theorem}[Product Formula]
\label{Lecture1:Theorem1.9}
\uses{Lecture1:Theorem1.8, ext:rationals}
[Theorem: Product Formula]
\end{theorem}

\begin{remark}[Section 1.4: Discrete valuations]
\label{Lecture1:Discussion_after_1.9}
\uses{Lecture1:Theorem1.9}
\discussion{Lecture1:Discussion_after_1.9}
[Discussion: Section 1.4: Discrete valuations]
\end{remark}

\begin{definition}[Valuation, value group, valuation ring, DVR]
\label{Lecture1:Definition1.10}
\uses{Lecture1:Discussion_after_1.9, ext:field, ext:ring, ext:integral-domain, ext:ideal-theory, ext:real-numbers, ext:fraction-field, ext:group-homomorphism, ext:discrete-subgroups-of-R}
[Definition: Valuation, value group, valuation ring, DVR]
\end{definition}

\begin{definition}[Valuation ring (general definition)]
\label{Lecture1:Definition1.11}
\uses{Lecture1:Definition1.10, ext:ring, ext:integral-domain, ext:ideal-theory, ext:fraction-field, ext:pid, ext:ufd}
[Definition: Valuation ring (general definition)]
\end{definition}

\begin{definition}[Local ring]
\label{Lecture1:Definition1.12}
\uses{Lecture1:Definition1.11, ext:ring, ext:ideal-theory}
[Definition: Local ring]
\end{definition}

\begin{definition}[Residue field]
\label{Lecture1:Definition1.13}
\uses{Lecture1:Definition1.12, ext:ideal-theory}
[Definition: Residue field]
\end{definition}

\begin{remark}[Determining the valuation from a DVR]
\label{Lecture1:Discussion_after_1.13}
\uses{Lecture1:Definition1.13, ext:ideal-theory, ext:fraction-field, ext:noetherian-ring, ext:krull-dimension, ext:regular-local-ring, ext:vector-space}
\discussion{Lecture1:Discussion_after_1.13}
[Discussion: Determining the valuation from a DVR]
\end{remark}

\begin{example}[p-adic valuation ring]
\label{Lecture1:Example1.14}
\uses{Lecture1:Discussion_after_1.13, ext:ideal-theory, ext:integers, ext:primes-and-divisibility, ext:finite-fields, ext:localization}
[Example: p-adic valuation ring]
\end{example}

\begin{example}[Laurent series valuation]
\label{Lecture1:Example1.15}
\uses{Lecture1:Example1.14, ext:field, ext:power-series-and-laurent-series}
[Example: Laurent series valuation]
\end{example}

\begin{remark}[Section 1.5: Properties of DVRs]
\label{Lecture1:Discussion_after_1.15}
\uses{Lecture1:Example1.15, ext:ideal-theory, ext:fraction-field}
\discussion{Lecture1:Discussion_after_1.15}
[Discussion: Section 1.5: Properties of DVRs]
\end{remark}

\begin{theorem}[Equivalent characterizations of DVRs]
\label{Lecture1:Theorem1.16}
\uses{Lecture1:Discussion_after_1.15, ext:integral-domain, ext:ideal-theory, ext:pid, ext:noetherian-ring, ext:krull-dimension, ext:localization, ext:regular-local-ring, ext:dvr-characterizations-proof}
[Theorem: Equivalent characterizations of DVRs]
\end{theorem}

\begin{remark}[Section 1.6: Integral extensions]
\label{Lecture1:Discussion_after_1.16}
\uses{Lecture1:Theorem1.16}
\discussion{Lecture1:Discussion_after_1.16}
[Discussion: Section 1.6: Integral extensions]
\end{remark}

\begin{definition}[Integral element and integral extension]
\label{Lecture1:Definition1.17}
\uses{Lecture1:Discussion_after_1.16, ext:ring, ext:polynomial-ring}
[Definition: Integral element and integral extension]
\end{definition}

\begin{proposition}[Sum and product of integral elements are integral]
\label{Lecture1:Proposition1.18}
\uses{Lecture1:Definition1.17, ext:ring, ext:polynomial-ring, ext:algebraic-closure, ext:symmetric-polynomials}
[Proposition: Sum and product of integral elements are integral]
\end{proposition}

\begin{definition}[Integral closure and normalization]
\label{Lecture1:Definition1.19}
\uses{Lecture1:Proposition1.18, ext:ring, ext:polynomial-ring, ext:fraction-field}
[Definition: Integral closure and normalization]
\end{definition}

\begin{proposition}[Transitivity of integrality]
\label{Lecture1:Proposition1.20}
\uses{Lecture1:Definition1.19, ext:ring, ext:transitivity-of-integrality-proof}
[Proposition: Transitivity of integrality]
\end{proposition}

\begin{corollary}[Integral closure is integrally closed]
\label{Lecture1:Corollary1.21}
\uses{Lecture1:Proposition1.20, ext:ring}
[Corollary: Integral closure is integrally closed]
\end{corollary}

\begin{proposition}[Z is integrally closed]
\label{Lecture1:Proposition1.22}
\uses{Lecture1:Corollary1.21, ext:integers, ext:rationals, ext:primes-and-divisibility, ext:polynomial-ring, ext:rational-root-test}
[Proposition: Z is integrally closed]
\end{proposition}

\begin{corollary}[Every UFD is integrally closed]
\label{Lecture1:Corollary1.23}
\uses{Lecture1:Proposition1.22, ext:integers, ext:pid, ext:ufd, ext:rational-root-test}
[Corollary: Every UFD is integrally closed]
\end{corollary}

\begin{example}[Z[sqrt(5)] is not integrally closed]
\label{Lecture1:Example1.24}
\uses{Lecture1:Corollary1.23, ext:integers, ext:polynomial-ring, ext:ufd}
[Example: Z[sqrt(5)] is not integrally closed]
\end{example}

\begin{proposition}[Every valuation ring is integrally closed]
\label{Lecture1:Proposition1.25}
\uses{Lecture1:Example1.24, ext:polynomial-ring, ext:fraction-field}
[Proposition: Every valuation ring is integrally closed]
\end{proposition}

\begin{definition}[Number field and ring of integers]
\label{Lecture1:Definition1.26}
\uses{Lecture1:Proposition1.25, ext:field, ext:integers, ext:rationals, ext:field-extensions}
[Definition: Number field and ring of integers]
\end{definition}

\begin{remark}[Notation for ring of integers]
\label{Lecture1:Remark1.27}
\uses{Lecture1:Definition1.26, ext:z-module, ext:vector-space}
\discussion{Lecture1:Remark1.27}
[Remark: Notation for ring of integers]
\end{remark}

\begin{proposition}[Minimal polynomial of integral element lies in A[x]]
\label{Lecture1:Proposition1.28}
\uses{Lecture1:Remark1.27, ext:integral-domain, ext:polynomial-ring, ext:fraction-field, ext:algebraic-closure, ext:field-extensions}
[Proposition: Minimal polynomial of integral element lies in A[x]]
\end{proposition}

\begin{example}[(1+sqrt(7))/2 is not integral over Z]
\label{Lecture1:Example1.29}
\uses{Lecture1:Proposition1.28, ext:integers, ext:polynomial-ring}
[Example: (1+sqrt(7))/2 is not integral over Z]
\end{example}

\begin{remark}[References]
\label{Lecture1:References}
\uses{Lecture1:Example1.29}
\discussion{Lecture1:References}
[Bibliography: References]
\end{remark}

\begin{remark}[MIT OpenCourseWare course information]
\label{Lecture1:Backmatter}
\uses{Lecture1:References}
\discussion{Lecture1:Backmatter}
[Discussion: MIT OpenCourseWare course information]
\end{remark}
